% ============================================================================
% Dual Coefficient Mappings -- rewritten under the margin-2 convention.
% Drop-in replacement for the previous \section{Dual Coefficient Mappings}.
%
% Requires in the preamble: amsmath, amsthm, bbm, algorithm, algorithmic, and
% theorem environments named: theorem, lemma, corollary, proposition, remark,
% example. e.g.
%   \newtheorem{theorem}{Theorem}
%   \newtheorem{lemma}{Lemma}
%   \newtheorem{corollary}{Corollary}[theorem]
%   \newtheorem{proposition}{Proposition}
%   \theoremstyle{remark} \newtheorem{remark}{Remark} \newtheorem{example}{Example}
%
% Cross-reference notes: this version is self-contained; equations that
% previously lived elsewhere (eq:SVM_Lagrangian_Dual_Loss, eq:Wolfe_Dual_Loss,
% eq:Box_Constraints, eq:pair_hinge, eq:Struct_Pair_Margin) are restated here
% with new labels. Search for "TODO" to reconcile with the rest of the paper.
%
% Summary of changes from the previous draft:
%   1. Margin-2 pairwise hinge throughout: [2 - <w, x_i - x_j>]_+. This is the
%      convention under which the pairwise and balanced dual OBJECTIVES agree
%      exactly (previously the linear terms differed by a factor of 2). A
%      remark shows margin-2 with parameter C is standard RankSVM with C/2,
%      so nothing is lost.
%   2. The injectivity claim of old Theorem 3 is removed (it is false for
%      n+, n- >= 2: a "rectangle move" on rho preserves all row/column sums).
%      It is replaced by the observation that the dual objective factors
%      through lambda, which is what the algorithm actually needs.
%   3. The image of the pairwise dual domain is now characterized EXACTLY:
%      it is a capacitated transportation polytope, strictly smaller than
%      box-plus-bias-hyperplane, with (a) a convex-hull description matching
%      the Frank-Wolfe vertices and (b) a Gale-type inequality description
%      giving a polynomial-time membership test via max-flow.
%   4. Old Theorems 1-2 are demoted to background lemmas (the error-rate
%      structural SVM / classification SVM equivalence is essentially known;
%      cite Joachims 2005, 2006 and Lacoste-Julien et al. 2013). Their proofs
%      are cleaned: the layer-cake construction is stated precisely and the
%      non-injectivity argument is now a counterexample, not a "contrapositive".
%   5. The bias-constraint corollary is promoted into the main equivalence
%      theorem: the pairwise dual lands in the WITH-intercept balanced dual
%      feasible set automatically.
%   6. New sandwich/certificate corollary: the balanced solution lower-bounds
%      the pairwise dual optimum and any transportation-feasible point upper-
%      bounds it, which is the foundation of the successive-tightening
%      algorithm in the next section.
% ============================================================================

\section{Dual Coefficient Mappings}
\label{sec:dual_maps}

This section establishes the precise relationship between the pairwise
(bipartite-ranking) SVM and the balanced cost-sensitive classification SVM at
the level of their dual programs. The main results are: (i) under the natural
row/column-sum map, the pairwise dual \emph{objective} coincides exactly with
the balanced classification dual objective (Theorem~\ref{thm:equivalence});
(ii) the pairwise dual feasible set maps onto a \emph{capacitated
transportation polytope} strictly contained in the balanced dual feasible set,
with the intercept constraint of the classification SVM arising automatically
(Theorems~\ref{thm:equivalence} and~\ref{thm:image}); and (iii) membership in
this polytope is testable in polynomial time
(Corollary~\ref{cor:membership}). Consequently, bipartite ranking with the
pairwise hinge \emph{is} balanced cost-sensitive classification with
additional transportation constraints on the dual variables, and each
Frank--Wolfe vertex of the pairwise dual corresponds to a sample-weighted
classification problem whose weights count misordered opposite-class examples
(Corollary~\ref{cor:fw_bridge}). These results license the dual-variable-free
algorithms of the next section: we may optimize the pairwise dual entirely in
the space of univariate Lagrangian coefficients, successively tightening the
balanced domain toward the transportation polytope, and stopping early
whenever leaving the polytope improves the margin.

\subsection{Setup and conventions}
\label{sec:conventions}

Let $x_i$, $i = 1, \dots, n^+$ denote the positive examples and $x_j$,
$j = 1, \dots, n^-$ the negative examples, with $N = n^+ + n^-$,
$p = n^+/N$, $q = n^-/N$, and labels $t_k \in \{+1, -1\}$ indexed by
$k = 1, \dots, N$ (index $i$ ranges over positives, $j$ over negatives, $k$
over all examples, and $\pi = (i,j)$ over positive--negative pairs). Write
$z_\pi = x_i - x_j$ for the pair difference vectors,
$Q = t t^\top \odot X X^\top$ for the standard SVM Gram operator, and
\begin{equation}
\bar{C} \;=\; \frac{C}{N^2}
\label{eq:cbar}
\end{equation}
for the per-pair regularization weight.

\paragraph{Pairwise hinge, margin-2 convention.}
The pairwise (RankSVM) primal is taken with margin $2$:
\begin{equation}
\min_{w}\;\; \frac{1}{2}\|w\|^2
\;+\; \bar{C} \sum_{i=1}^{n^+} \sum_{j=1}^{n^-}
\bigl[\, 2 - \langle w, x_i - x_j \rangle \,\bigr]_+ .
\label{eq:pair_primal}
\end{equation}
Its Lagrangian dual, with one multiplier $\rho_\pi = \rho_{ij}$ per pair, is
\begin{equation}
\min_{\rho}\;\; \frac{1}{2} \rho^\top \widetilde{Q} \rho
\;-\; 2\,\mathbbm{1}^\top \rho
\qquad \text{s.t.} \quad 0 \le \rho_\pi \le \bar{C} \;\;\forall \pi,
\label{eq:pair_dual}
\end{equation}
where $\widetilde{Q}_{\pi\pi'} = \langle z_\pi, z_{\pi'} \rangle$ and
$w = \sum_\pi \rho_\pi z_\pi$ at optimality. There is no equality constraint
because the intercept cancels within each pair.

\begin{remark}[Margin~2 is standard RankSVM, rescaled]
\label{rem:margin2}
For any $w$, substituting $v = w/2$ in \eqref{eq:pair_primal} gives
\[
\frac{1}{2}\|w\|^2 + \bar{C}\sum_{\pi}\bigl[2 - \langle w, z_\pi\rangle\bigr]_+
= 4\left( \frac{1}{2}\|v\|^2
+ \frac{\bar{C}}{2}\sum_{\pi}\bigl[1 - \langle v, z_\pi\rangle\bigr]_+ \right).
\]
Hence the margin-2 problem with parameter $C$ has the same minimizer, up to
the scaling $w = 2v$, as the unit-margin RankSVM with parameter $C/2$. The
induced ranking function, the AUC, and the regularization path are unchanged;
only the parameterization of $C$ shifts. The margin-2 convention is the
natural one here because it is the margin produced by summing the two
univariate hinges of a pair:
$[1 - \langle w, x_i\rangle]_+ + [1 + \langle w, x_j\rangle]_+
\ge [2 - \langle w, x_i - x_j\rangle]_+$.
\end{remark}

\paragraph{Balanced classification.}
The balanced (class-weighted) binary SVM with intercept weights each positive
example by $\bar{C} n^-$ and each negative example by $\bar{C} n^+$:
\begin{equation}
\min_{w, b}\;\; \frac{1}{2}\|w\|^2
\;+\; \bar{C}\, n^- \sum_{i=1}^{n^+} \bigl[ 1 - (\langle w, x_i\rangle + b) \bigr]_+
\;+\; \bar{C}\, n^+ \sum_{j=1}^{n^-} \bigl[ 1 + (\langle w, x_j\rangle + b) \bigr]_+ .
\label{eq:bal_primal}
\end{equation}
(Equivalently, per-class weights $\tfrac{C}{N}q$ and $\tfrac{C}{N}p$: each
example is weighted by the empirical frequency of the \emph{opposite} class.)
Its dual, with one multiplier $\lambda_k$ per example, is
\begin{equation}
\min_{\lambda}\;\; D(\lambda) \;:=\; \frac{1}{2}\lambda^\top Q \lambda
\;-\; \mathbbm{1}^\top \lambda
\qquad \text{s.t.} \quad \lambda \in \mathcal{B},
\label{eq:bal_dual}
\end{equation}
\begin{equation}
\mathcal{B} \;=\; \Bigl\{ \lambda \in \mathbbm{R}^N_{\ge 0} \;:\;
\lambda_i \le \bar{C} n^- \;\forall i, \;\;
\lambda_j \le \bar{C} n^+ \;\forall j, \;\;
\textstyle\sum_k t_k \lambda_k = 0 \Bigr\},
\label{eq:bal_domain}
\end{equation}
with $w = \sum_k t_k \lambda_k x_k$ at optimality. The equality constraint in
\eqref{eq:bal_domain} is the usual consequence of the intercept $b$.

\subsection{Background: Frank--Wolfe and Lagrangian dual coefficients}
\label{sec:background}

The structural SVM and its (block-coordinate) Frank--Wolfe optimizers
% TODO: cite Joachims (2005, 2006); Lacoste-Julien, Jaggi, Schmidt, Pletscher (2013)
work with dual coefficients $\alpha$ indexed by \emph{joint labelings}
(vertices) rather than by examples. For a decomposable hinge, each vertex
$s \in \{0,1\}^N$ marks which examples are margin-violated, and the linear
operator $\mathcal{S}$, whose columns are the vertices, maps vertex
coefficients to example coefficients:
\begin{equation}
\lambda \;=\; \mathcal{S}\alpha,
\qquad
\alpha \in \Delta_C := \Bigl\{ \alpha \ge 0 : \textstyle\sum_s \alpha_s = C \Bigr\}.
\label{eq:wolfe_map}
\end{equation}
The two background lemmas below record that nothing is lost by tracking
$\lambda$ instead of $\alpha$; they are essentially the known equivalence
between the error-rate structural SVM and the classification SVM
% TODO: cite Joachims (2006), Theorem 1
and are stated in the form used later.

\begin{lemma}[Per-example blocks: bijection]
\label{lem:block_bijection}
In the block-separable formulation (one block per example, block mass $C_k$
equal to that example's primal weight), each block has two vertices --
violated and satisfied -- with coefficients $\alpha_{k1} + \alpha_{k0} = C_k$.
The map $\lambda_k = \alpha_{k1}$ is a bijection between the per-block
Frank--Wolfe domain and the Lagrangian box $\lambda_k \in [0, C_k]$, applied
independently across blocks.
\end{lemma}

\begin{proof}
Immediate: $\alpha_{k1} \in [0, C_k]$ determines $\alpha_{k0} = C_k -
\alpha_{k1}$ uniquely, and every $\lambda_k \in [0, C_k]$ is attained.
\end{proof}

\begin{lemma}[Holistic formulation: surjection onto the box]
\label{lem:holistic_surjection}
Let $\mathcal{S} \in \{0,1\}^{N \times 2^N}$ contain all binary vertices as
columns. The map $\alpha \mapsto \mathcal{S}\alpha$ sends $\Delta_C$
\emph{onto} $[0, C]^N$, using at most $N$ vertices with nested supports; it is
not injective for $N \ge 2$.
\end{lemma}

\begin{proof}
(\emph{Surjectivity: layer-cake construction.}) Fix $\lambda \in [0,C]^N$ and
let $0 < v_1 < v_2 < \cdots < v_r$ be its distinct positive values, with
$v_0 = 0$. Define nested vertices $s_m = \mathbbm{1}\{\lambda \ge v_m\}$ and
coefficients $\alpha_{s_m} = v_m - v_{m-1} > 0$ for $m = 1, \dots, r$. For each
coordinate $k$, $\sum_m \alpha_{s_m} (s_m)_k = \sum_{m : v_m \le \lambda_k}
(v_m - v_{m-1}) = \lambda_k$, so $\mathcal{S}\alpha = \lambda$. The
coefficients sum to $v_r = \max_k \lambda_k \le C$; assign the remainder
$C - v_r$ to the zero vertex $s_0 = \mathbf{0}$. This uses $r \le N$ nonzero
vertices plus the zero vertex, all with nested supports.

(\emph{Non-injectivity.}) For $N = 2$, the distinct points
$\alpha = \tfrac{C}{2}\bigl(e_{(1,1)} + e_{(0,0)}\bigr)$ and
$\widetilde{\alpha} = \tfrac{C}{2}\bigl(e_{(1,0)} + e_{(0,1)}\bigr)$
both map to $\lambda = (\tfrac{C}{2}, \tfrac{C}{2})$.
\end{proof}

For the pairwise loss, the vertices are subsets $s$ of positive--negative
pairs, and the corresponding column of $\mathcal{S}$ records, for each
example, \emph{how many} pairs in $s$ it participates in:
\begin{equation}
(\mathcal{S} e_s)_i \;=\; \#\{\, j : (i,j) \in s \,\}, \qquad
(\mathcal{S} e_s)_j \;=\; \#\{\, i : (i,j) \in s \,\}.
\label{eq:pair_counts}
\end{equation}
Under the margin-2 convention the structural loss of vertex $s$ is
$\Delta_s = 2|s|$ (each misordered pair contributes its margin), and since
each pair touches exactly two examples,
$\mathbbm{1}^\top \mathcal{S} e_s = 2|s| = \Delta_s$. Summing over vertices:
\begin{equation}
\Delta^\top \alpha \;=\; \mathbbm{1}^\top \mathcal{S} \alpha
\;=\; \mathbbm{1}^\top \lambda .
\label{eq:margin_identity}
\end{equation}
This is the pairwise analogue of the classification identity and is the
bookkeeping that makes Theorem~\ref{thm:equivalence} exact. (Under a
unit-margin convention the two sides differ by a factor of $2$; see
Remark~\ref{rem:margin2}.)

\subsection{The pairwise dual is a constrained balanced dual}
\label{sec:equivalence}

Define the row/column-sum map $\Phi$ from pair coefficients to example
coefficients:
\begin{equation}
\Phi(\rho) \;=\; \lambda, \qquad
\lambda_i = \sum_{j=1}^{n^-} \rho_{ij}, \qquad
\lambda_j = \sum_{i=1}^{n^+} \rho_{ij}.
\label{eq:phi_map}
\end{equation}

\begin{theorem}[Dual equivalence]
\label{thm:equivalence}
Let $\mathcal{R} = [0, \bar{C}]^{\,n^+ \times n^-}$ be the pairwise dual
feasible set \eqref{eq:pair_dual} and let $\lambda = \Phi(\rho)$. Then:
\begin{enumerate}
\item[(i)] \emph{(Weight vectors agree.)}
$\displaystyle \sum_\pi \rho_\pi z_\pi = \sum_k t_k \lambda_k x_k$.
\item[(ii)] \emph{(Objectives agree.)}
$\displaystyle \frac{1}{2}\rho^\top \widetilde{Q} \rho - 2\,\mathbbm{1}^\top\rho
= \frac{1}{2}\lambda^\top Q \lambda - \mathbbm{1}^\top\lambda = D(\lambda)$.
\item[(iii)] \emph{(Feasibility is preserved, intercept included.)}
$\Phi(\mathcal{R}) \subseteq \mathcal{B}$, where $\mathcal{B}$ is the
balanced feasible set \eqref{eq:bal_domain}. In particular the equality
constraint $\sum_k t_k \lambda_k = 0$ holds automatically for every
$\lambda \in \Phi(\mathcal{R})$.
\end{enumerate}
Consequently, with $\mathcal{T} := \Phi(\mathcal{R})$,
\begin{equation}
\min_{\rho \in \mathcal{R}}\;
\frac{1}{2}\rho^\top \widetilde{Q} \rho - 2\,\mathbbm{1}^\top\rho
\;\;=\;\;
\min_{\lambda \in \mathcal{T}}\; D(\lambda),
\label{eq:reduced_problem}
\end{equation}
i.e.\ the pairwise dual is exactly the balanced classification dual restricted
to the subdomain $\mathcal{T} \subseteq \mathcal{B}$.
\end{theorem}

\begin{proof}
(i) Splitting each pair difference,
\[
\sum_\pi \rho_\pi z_\pi
= \sum_{i}\sum_{j} \rho_{ij}(x_i - x_j)
= \sum_i \Bigl(\sum_j \rho_{ij}\Bigr) x_i - \sum_j \Bigl(\sum_i \rho_{ij}\Bigr) x_j
= \sum_i \lambda_i x_i - \sum_j \lambda_j x_j
= \sum_k t_k \lambda_k x_k .
\]

(ii) The quadratic terms both equal $\|w\|^2$ for the common $w$ of part (i):
$\rho^\top \widetilde{Q}\rho = \|\sum_\pi \rho_\pi z_\pi\|^2 =
\|\sum_k t_k\lambda_k x_k\|^2 = \lambda^\top Q \lambda$. For the linear terms,
each unit of $\rho_{ij}$ contributes once to $\lambda_i$ and once to
$\lambda_j$, so
\[
\mathbbm{1}^\top \lambda
= \sum_i \lambda_i + \sum_j \lambda_j
= 2 \sum_\pi \rho_\pi
= 2\,\mathbbm{1}^\top \rho .
\]

(iii) Row and column sums of a matrix with entries in $[0, \bar{C}]$ obey
$\lambda_i = \sum_j \rho_{ij} \le \bar{C} n^-$ and
$\lambda_j = \sum_i \rho_{ij} \le \bar{C} n^+$, matching the balanced box in
\eqref{eq:bal_domain}. Moreover $\sum_i \lambda_i = \sum_{ij} \rho_{ij} =
\sum_j \lambda_j$, i.e.\ $\sum_k t_k \lambda_k = 0$: the total mass on
positives equals the total mass on negatives because every pair contributes
equally to both sides. Thus $\Phi(\rho) \in \mathcal{B}$.

Finally, \eqref{eq:reduced_problem} follows because by (ii) the objective is
constant on the fibers of $\Phi$, so minimizing over $\mathcal{R}$ and
minimizing $D$ over the image $\mathcal{T}$ are the same problem.
\end{proof}

\begin{remark}[$\Phi$ is not injective, and this does not matter]
\label{rem:noninjective}
For $n^+, n^- \ge 2$ the map $\Phi$ cannot be injective: it maps
$n^+ n^-$ dimensions to $n^+ + n^-$ dimensions, and explicitly, the
\emph{rectangle move} replacing
$(\rho_{ij},\, \rho_{i'j'},\, \rho_{ij'},\, \rho_{i'j})$ by
$(\rho_{ij} + \varepsilon,\, \rho_{i'j'} + \varepsilon,\,
\rho_{ij'} - \varepsilon,\, \rho_{i'j} - \varepsilon)$
preserves every row and column sum. However, by
Theorem~\ref{thm:equivalence}(ii) the dual objective depends on $\rho$ only
through $\lambda = \Phi(\rho)$, so the reduced problem
\eqref{eq:reduced_problem} loses nothing: any minimizer $\lambda^\star$ of $D$
over $\mathcal{T}$ pulls back to a (non-unique) optimal $\rho$, and the primal
solution $w$ is recovered from $\lambda^\star$ alone.
\end{remark}

\subsection{The image is a capacitated transportation polytope}
\label{sec:image}

Theorem~\ref{thm:equivalence} reduces the pairwise dual to minimizing the
balanced dual objective over $\mathcal{T} = \Phi(\mathcal{R})$. We now
characterize $\mathcal{T}$ exactly. Write $\Lambda := \sum_i \lambda_i =
\sum_j \lambda_j$ for the common one-sided mass of a point satisfying the
equality constraint, and let $L^+(a)$ (resp.\ $L^-(b)$) denote the sum of the
$a$ largest positive-example coefficients (resp.\ $b$ largest
negative-example coefficients).

\begin{theorem}[Exact image]
\label{thm:image}
$\mathcal{T}$ admits the following two descriptions.
\begin{enumerate}
\item[(i)] \emph{(Vertex description.)} $\mathcal{T}$ is the convex hull of
the scaled misordering-count vectors:
\[
\mathcal{T} \;=\; \operatorname{conv}
\bigl\{\, \bar{C}\, \mathcal{S} e_s \;:\; s \subseteq \{1,\dots,n^+\} \times
\{1,\dots,n^-\} \,\bigr\},
\]
with $\mathcal{S}e_s$ the per-example pair counts of \eqref{eq:pair_counts}.
\item[(ii)] \emph{(Inequality description.)} $\lambda \in \mathcal{T}$ if and
only if $\lambda \ge 0$, $\sum_i \lambda_i = \sum_j \lambda_j = \Lambda$, and
\begin{equation}
L^+(a) + L^-(b) \;\le\; \Lambda + \bar{C}\, a b
\qquad \text{for all } 1 \le a \le n^+, \; 1 \le b \le n^- .
\label{eq:gale}
\end{equation}
\end{enumerate}
Moreover $\mathcal{T} \subsetneq \mathcal{B}$ whenever $n^+, n^- \ge 2$: the
box and equality constraints of $\mathcal{B}$ do not imply \eqref{eq:gale}.
\end{theorem}

\begin{proof}
(i) The feasible set $\mathcal{R} = [0,\bar{C}]^{n^+ \times n^-}$ is a
polytope whose vertices are $\bar{C}\,\mathbbm{1}_s$ for pair subsets $s$
(entries equal to $\bar{C}$ on $s$, zero elsewhere). The image of a polytope
under a linear map is the convex hull of the images of its vertices, and
$\Phi(\bar{C}\,\mathbbm{1}_s) = \bar{C}\,\mathcal{S}e_s$ by
\eqref{eq:pair_counts}.

(ii) The question is feasibility of a transportation problem: does there
exist $\rho \in [0,\bar{C}]^{n^+ \times n^-}$ with prescribed row sums
$\lambda_i$ and column sums $\lambda_j$? Model it as a flow network: source
$\sigma \to$ row $i$ with capacity $\lambda_i$; row $i \to$ column $j$ with
capacity $\bar{C}$; column $j \to$ sink $\tau$ with capacity $\lambda_j$.
Feasibility is equivalent to a maximum flow of value $\Lambda$, which by
max-flow/min-cut holds iff every $\sigma$--$\tau$ cut has capacity at least
$\Lambda$. A cut is determined by the set $S$ of rows and the set $T'$ of
columns on the source side; writing $T$ for the complement of $T'$, its
capacity is
$\sum_{i \notin S}\lambda_i + \bar{C}\,|S|\,|T| + \sum_{j \in
T'}\lambda_j$. Requiring this to be at least $\Lambda = \sum_i \lambda_i$
yields
\[
\sum_{i \in S} \lambda_i \;\le\; \bar{C}\,|S|\,|T| + \sum_{j \notin T} \lambda_j
\qquad \forall\, S, T .
\]
For fixed cardinalities $a = |S|$, $b = |T|$, the left side is maximized by
the $a$ largest $\lambda_i$ and the right side minimized by placing the $b$
largest $\lambda_j$ inside $T$; substituting
$\sum_{j \notin T}\lambda_j = \Lambda - L^-(b)$ gives \eqref{eq:gale}. Cases
$a = 0$ or $b = 0$ are implied by nonnegativity and the equal-sum condition.

(\emph{Strictness.}) Example~\ref{ex:strict} below exhibits, for
$n^+ = n^- = 2$, a point of $\mathcal{B}$ violating \eqref{eq:gale}.
\end{proof}

\begin{example}[$\mathcal{T}$ is strictly smaller than
box $\cap$ hyperplane]
\label{ex:strict}
Let $n^+ = n^- = 2$ and consider
$\lambda^+ = (2\bar{C},\, 0)$, $\lambda^- = (2\bar{C},\, 0)$. All box
constraints of $\mathcal{B}$ hold ($2\bar{C} = \bar{C}n^-$) and the sums are
equal, so $\lambda \in \mathcal{B}$. But with $a = b = 1$,
$L^+(1) + L^-(1) = 4\bar{C} > \Lambda + \bar{C} = 3\bar{C}$, violating
\eqref{eq:gale}: the single positive example carrying mass $2\bar{C}$ would
have to route more than the cell capacity $\bar{C}$ through the single
negative example carrying mass. Concretely, no $\rho \in [0,\bar{C}]^{2\times
2}$ has row sums $(2\bar{C}, 0)$ and column sums $(2\bar{C}, 0)$.
\end{example}

\begin{corollary}[Polynomial-time membership and projection oracles]
\label{cor:membership}
Whether $\lambda \in \mathcal{T}$ can be decided by a single maximum-flow
computation on the bipartite network of Theorem~\ref{thm:image}(ii) (with
$n^+ + n^- + 2$ nodes), or by checking the $n^+ n^-$ sorted prefix conditions
\eqref{eq:gale} after two sorts. A witness $\rho$ with $\Phi(\rho) = \lambda$
is recovered from the flow itself. Consequently, the successive-tightening
algorithms of Section~\ref{sec:algorithms} have an efficient certificate for
``the current iterate is (nearly) pairwise-feasible,'' and Euclidean
projection onto $\mathcal{T}$ is a quadratic program over the same flow
polytope.
% TODO: fix \ref{sec:algorithms} to the actual label of the Algorithms section.
\end{corollary}

\subsection{Consequences}
\label{sec:consequences}

\begin{corollary}[Sandwich certificate for the pairwise optimum]
\label{cor:sandwich}
Let $\lambda^\star_{\mathrm{bal}} \in \arg\min_{\lambda \in \mathcal{B}}
D(\lambda)$ and $\lambda^\star_{\mathrm{pair}} \in \arg\min_{\lambda \in
\mathcal{T}} D(\lambda)$. Then for every $\lambda \in \mathcal{T}$,
\begin{equation}
D(\lambda^\star_{\mathrm{bal}})
\;\le\; D(\lambda^\star_{\mathrm{pair}})
\;\le\; D(\lambda),
\label{eq:sandwich}
\end{equation}
so any pairwise-feasible iterate $\lambda \in \mathcal{T}$ certifies its own
suboptimality on the pairwise problem:
$D(\lambda) - D(\lambda^\star_{\mathrm{pair}})
\le D(\lambda) - D(\lambda^\star_{\mathrm{bal}})$.
In particular, if the balanced solution happens to satisfy
$\lambda^\star_{\mathrm{bal}} \in \mathcal{T}$, it \emph{is} a pairwise
(RankSVM) dual solution.
\end{corollary}

\begin{proof}
Immediate from $\mathcal{T} \subseteq \mathcal{B}$
(Theorem~\ref{thm:equivalence}(iii)) and the shared objective $D$
(Theorem~\ref{thm:equivalence}(ii)).
\end{proof}

\begin{corollary}[Frank--Wolfe vertices are weighted classification problems]
\label{cor:fw_bridge}
Each Frank--Wolfe vertex $s$ of the pairwise dual corresponds under
$\lambda_s = \bar{C}\,\mathcal{S}e_s$ to the vector of \emph{misordering
counts}: example $i$ receives weight proportional to the number of negatives
in $s$ paired above it, and example $j$ to the number of positives paired
below it. By Theorem~\ref{thm:image}(i), $\lambda_s \in \mathcal{T}$, and
$\mathcal{T}$ is exactly the convex hull of these vertices; by
\eqref{eq:margin_identity} the Frank--Wolfe linearization
$\langle \nabla, \alpha \rangle$-objective transfers to $\lambda$-space with
the same value. Hence running Frank--Wolfe on the pairwise dual is running
Frank--Wolfe on $\min_{\lambda \in \mathcal{T}} D(\lambda)$, and each
linear-minimization step selects the count-weighted univariate problem used
by the algorithms of the next section. At the first iterate ($w = 0$) every
pair is violated, so the counts are $n^-$ for every positive and $n^+$ for
every negative: \emph{the balanced classification problem
\eqref{eq:bal_primal} is precisely the first Frank--Wolfe subproblem}.
\end{corollary}

\begin{remark}[Relaxed trust regions]
\label{rem:relaxation}
The trust-region subproblems of Section~\ref{sec:algorithms} optimize $D$
over boxes $[0, \bar{C}\omega]$ with $\omega = \mathcal{S}e_s$ the current
counts. The vertex $\bar{C}\omega$ itself lies in $\mathcal{T}$, but the box
is \emph{not} contained in $\mathcal{T}$ (coordinate-wise shrinking breaks the
equal-mass constraint, and \eqref{eq:gale} may fail). The subproblems are
therefore relaxations: their solutions can leave $\mathcal{T}$, and by
Corollary~\ref{cor:sandwich} they lower-bound the pairwise optimum whenever
they leave it. The algorithms exploit this deliberately -- an iterate is
pulled back toward $\mathcal{T}$ only when doing so does not worsen the
margin -- and Corollary~\ref{cor:membership} supplies the test.
% TODO: fix \ref{sec:algorithms} to the actual label of the Algorithms section.
\end{remark}

% ----------------------------------------------------------------------------
% Figure: update the old Coeff_Maps.png caption to match the new statement.
% Suggested replacement:
%
% \begin{figure}
%     \centering
%     \includegraphics[scale=.8]{Figures/Coeff_Maps.png}
%     \caption{Dual coefficient domains. The pairwise Frank--Wolfe simplex
%     maps (via the misordering-count operator $\mathcal{S}$) onto the
%     capacitated transportation polytope $\mathcal{T}$, which sits strictly
%     inside the balanced classification dual feasible set $\mathcal{B}$
%     (box $\cap$ intercept hyperplane). The dual objective $D$ is identical
%     on $\mathcal{T}$ for both problems (Theorem~\ref{thm:equivalence}).}
%     \label{fig:coeff_maps}
% \end{figure}
% ----------------------------------------------------------------------------
